\subsection{幂零自同态的线性李代数}

引理 1. 设 \(\mathfrak{n}\) 是 \(\mathfrak{gl}(\mathrm{V})\) 的一个由幂零自同态组成的李子代数，\(\mathrm{N}\) 是子群 \(\exp \mathfrak{n}\)（\(\mathbf{GL}(\mathrm{V})\) 的子群）（§3，第 1 节，引理 1）。

\begin{enumerate}
\def\labelenumi{(\roman{enumi})}
\item
  设 \(\rho\) 是 \(\mathfrak{n}\) 在 W 上的一个有限维线性表示，使得 \(\rho(\mathfrak{n})\) 的元素都是幂零的，W' 是 W 的一个在 \(\rho\) 下稳定的向量子空间，\(\rho_1\) 与 \(\rho_2\) 是由 W' 定义的 \(\rho\) 的子表示与商表示，\(\pi, \pi_1, \pi_2\) 是与 \(\rho, \rho_1, \rho_2\) 相容的 N 的表示（§3，第 1 节）。则 \(\pi_1, \pi_2\) 是由 W' 定义的 \(\pi\) 的子表示与商表示。
\item
  设 \(\rho_1, \rho_2\) 是 \(\mathfrak{n}\) 的有限维线性表示，使得 \(\rho_1(\mathfrak{n})\) 与 \(\rho_2(\mathfrak{n})\) 的元素都是幂零的，\(\pi_1, \pi_2\) 是与 \(\rho_1, \rho_2\) 相容的 N 的表示。则 \(\pi_1 \otimes \pi_2\) 是与 \(\rho_1 \otimes \rho_2\) 相容的 N 的表示。
\item
  设 \(\rho_{1},\rho_{2}\) 是 \(\mathfrak{n}\) 在向量空间 \(V_{1},V_{2}\) 上的有限维线性表示，使得 \(\rho_1(\mathfrak{n})\) 与 \(\rho_{2}(\mathfrak{n})\) 的元素都是幂零的，\(\rho\) 是 \(\mathfrak{n}\) 在 \(\mathrm{Hom}(V_1,V_2)\) 上由 \(\rho_{1},\rho_{2}\) 确定的表示。设 \(\pi_1,\pi_2\) 是与 \(\rho_{1},\rho_{2}\) 相容的 N 的表示，\(\pi\) 是 N 在 \(\mathrm{Hom}(\mathrm{V}_1,\mathrm{V}_2)\) 上由 \(\pi_1,\pi_2\) 确定的表示。则 \(\pi\) 是与 \(\rho\) 相容的 N 的表示。
\end{enumerate}

断言 (i) 是明显的。设 \(\rho_1, \rho_2, \pi_1, \pi_2\) 如 (ii) 中所述。若 \(x \in \mathfrak{n}\)，则由于 \(\rho_1(x) \otimes 1\) 与 \(1 \otimes \rho_2(x)\) 交换，我们有

\[
\begin{array}{l} \exp (\rho_ {1} (x) \otimes 1 + 1 \otimes \rho_ {2} (x)) = \exp (\rho_ {1} (x) \otimes 1). \exp (1 \otimes \rho_ {2} (x)) \\ \qquad = (\exp \rho_ {1} (x)) \otimes 1. 1 \otimes (\exp \rho_ {2} (x)) \\ \qquad = (\exp \rho_ {1} (x)) \otimes (\exp \rho_ {2} (x)) \\ \qquad = \pi_ {1} (\exp x) \otimes \pi_ {2} (\exp x) \\ \qquad = (\pi_ {1} \otimes \pi_ {2}) (\exp x), \end{array}
\]

故得 (ii)。设 \(\rho_1, \rho_2, \rho, \pi_1, \pi_2, \pi, V_1, V_2\) 如 (iii) 中所述。若 \(v_1 \in \mathrm{End}V_1\) 且 \(v_2 \in \mathrm{End}V_2\)，用 \(R_{v_1}\) 与 \(L_{v_2}\) 表示映射 \(u \mapsto uv_1\) 与 \(u \mapsto v_2u\)（从 \(\mathrm{Hom}(V_1, V_2)\) 到自身的映射）；这些映射交换，且 \(\rho(x)u = (L_{\rho_2(x)} - R_{\rho_1(x)})u\)，于是

\[
\begin{array}{r l} & {\exp \rho (x). u = \exp \mathrm{L} _ {\rho_ {2} (x)}. \exp \mathrm{R} _ {- \rho_ {1} (x)}. u} \\ & {\qquad = \mathrm{L} _ {\exp \rho_ {2} (x)}. \mathrm{R} _ {\exp (- \rho_ {1} (x))}. u} \\ & {\qquad = \mathrm{L} _ {\pi_ {2} (\exp x)}. \mathrm{R} _ {\pi_ {1} (\exp (- x))}. u} \\ & {\qquad = \pi (\exp x). u,} \end{array}
\]

故得 (iii)。

引理 2 \(^{2}\). (i) 设 W 是 V 的一个维数为 d 的向量子空间，D 是直线 \(\bigwedge^{d}\) W \(\subset \bigwedge^{d}\) V，\(\theta\) 是 \(\mathfrak{gl}(V)\) 在 \(\bigwedge V\) 上的典则表示（第三章，附录）。设 \(x \in \mathfrak{gl}(V)\)。则 \(x(W) \subset W\) 当且仅当 \(\theta(x)(D) \subset D\)。

\begin{enumerate}
\def\labelenumi{(\roman{enumi})}
\setcounter{enumi}{1}
\tightlist
\item
  设 \((e_{1},\ldots,e_{n})\) 是 \(k^{n}\) 的典则基，\(\theta\) 是 \(\mathfrak{gl}(n,k)\) 在 \(\bigwedge(k^{n})\) 上的典则表示，\(x\in\mathfrak{gl}(n,k)\)。则 \(x\in\mathfrak{n}(n,k)\) 当且仅当
\end{enumerate}

\[
\theta (x) (e _ {n - d + 1} \wedge \dots \wedge e _ {n}) = 0
\]

对所有 \(1 \leq d \leq n\) 成立。

\footnotetext[2]{在这个引理中，\(k\) 可以是任意（交换）域。}

\begin{enumerate}
\def\labelenumi{(\roman{enumi})}
\tightlist
\item
  若 \(x(\mathrm{W}) \subset \mathrm{W}\)，则显然 \(\theta(x)\mathrm{D} \subset \mathrm{D}\)。反之，设 \(\theta(x)\mathrm{D} \subset \mathrm{D}\)。设 \(u\) 是 \(\mathrm{D}\) 的一个非零元素，并设 \(y \in \mathrm{W}\)。则 \(y \wedge u = 0\)。由于 \(\theta(x)\) 是 \(\bigwedge \mathrm{V}\) 的一个导子，这蕴含
\end{enumerate}

\[
\theta (x) y \wedge u + y \wedge \theta (x) u = 0.
\]

而现在 \(\theta(x)u\in ku\)，故 \(y\wedge\theta(x)u=0\)，从而 \(\theta(x)y\wedge u=0\)。由《代数》第三章 §7 第 9 节命题 13，这蕴含 \(\theta(x)y\in W\)，即 \(x(y)\in W\)，这就证明 \(x(W)\subset W\)。

\begin{enumerate}
\def\labelenumi{(\roman{enumi})}
\setcounter{enumi}{1}
\tightlist
\item
  对 \(x \in \mathfrak{n}(n, k)\) 而言，(ii) 中所述的条件显然是必要的。设它成立。由 (i)，\(x\) 使
\end{enumerate}

\[
k e _ {n - d + 1} + \dots + k e _ {n}
\]

保持稳定，又由于这对 \(d = 1, \ldots, n\) 都成立，\(x\) 是下三角的。置

\[
x = (x _ {i j}) _ {1 \leq i, j \leq n}.
\]

我们有 \(0 = x(e_n) = x_{nn}e_n\)，故 \(x_{nn} = 0\)。设 \(i < n\)，并设已证明 \(x_{jj} = 0\) 对 \(j > i\) 成立。则

\[
0 = \theta (x) \left(e _ {i} \wedge e _ {i + 1} \wedge \dots \wedge e _ {n}\right) = x _ {i i} \left(e _ {i} \wedge e _ {i + 1} \wedge \dots \wedge e _ {n}\right),
\]

故 \(x_{ii} = 0\)。于是 \(x \in \mathfrak{n}(n,k)\)。

命题 8. 设 \(\mathfrak{n}\) 是 \(\mathfrak{gl}(\mathrm{V})\) 的一个由幂零元素组成的李子代数，\(\mathfrak{q}\) 是 \(\mathfrak{n}\) 在 \(\mathfrak{gl}(\mathrm{V})\) 中的正规化子。存在一个有限维向量空间 E、一个表示 \(\rho\)（\(\mathfrak{gl}(\mathrm{V})\) 在 E 上的表示）以及 E 的一个向量子空间 F，满足下列条件：

\begin{enumerate}
\def\labelenumi{(\roman{enumi})}
\item
  在 \(\rho\) 下，\(V\) 的位似的像是可对角化的；
\item
  F 在 \(\rho(\mathfrak{q})\) 下稳定；
\item
  \(\mathfrak{n}\) 是那些 \(x\in \mathfrak{gl}(\mathrm{V})\)（满足 \(\rho (x)(\mathrm{F}) = 0\)）所成的集
\end{enumerate}

设 \(n = \dim V\)。由 Engel 定理，可以把 \(V\) 与 \(k^n\) 等同起来，使得 \(\mathfrak{n} \subset \mathfrak{n}(n,k)\)。设 \(P\) 是 \(\mathfrak{gl}(n,k)\) 上的多项式函数所成的代数。对 \(i = 0,1,\ldots\)，设 \(P_i\) 是 \(P\) 中 \(i\) 次齐次元素所成的集。设 \(N = \exp \mathfrak{n}\)，它是严格下三角群 \(T\) 的一个子群。设 \(J\) 是 \(P\) 中在 \(N\) 上为零的元素所成的集；这是 \(P\) 中的一个理想。设 \(N_J\) 是那些 \(x \in \mathfrak{gl}(n,k)\)（使得 \(p(x) = 0\) 对所有 \(p \in J\) 成立）所成的集。则 \(N \subset N_J\)。反之，设 \(x \in N_J\)。用 \(p_{ij}\) 表示给出 \(\mathfrak{gl}(n,k)\) 的一个元素的矩阵项的多项式函数。理想 \(J\) 含有诸 \(p_{ij}\)（对 \(i < j\)）与诸 \(p_{ii} - 1\)；因而 \(x \in T\)。另一方面，若 \(u\) 是 \(\mathfrak{gl}(n,k)\) 上的一个在 \(\mathfrak{n}\) 上为零的线性形式，则存在 \(p_u \in P\)，使得 \(p_u(z) = u(\log z)\) 对所有 \(z \in T\) 成立（\(\S 3\)，第 1 节，引理 1 (i)）；我们有 \(p_u \in J\)，故 \(u(\log x) = 0\)。由此推出 \(\log x\) 属于 \(\mathfrak{n}\)，于是 \(x \in N\)，这就证明了 \(N = N_J\)。

对所有 \(p \in P\) 与 \(g \in \mathbf{GL}_{n}(k)\)，设 \(\lambda(g)p\) 是函数 \(x \mapsto p(g^{-1}x)\)（\(\mathfrak{gl}(n,k)\) 上的函数）；则 \(\lambda(g)p \in \mathrm{P}\)，\(\lambda(g)\) 是代数 P 的一个自同构，且 \(\lambda\) 是 \(\mathbf{GL}_{n}(k)\) 在 P 上的一个表示，它使每个 \(P_{i}\) 保持稳定。我们证明

\[
\mathrm{N} = \{x \in \mathbf {G L} _ {n} (k) \mid \lambda (x) \mathrm{J} = \mathrm{J} \}.\tag{1}
\]

若 \(x \in \mathrm{N}, p \in \mathrm{J}, y \in \mathrm{N}\)，则 \((\lambda(x)p)(y) = p(x^{-1}y) = 0\)，因为 \(x^{-1}y \in \mathrm{N}\)；于是 \(\lambda(x)p \in \mathrm{J}\)，故 \(\lambda(x)\mathrm{J} = \mathrm{J}\)。设 \(x \in \mathbf{GL}_n(k)\) 满足 \(\lambda(x)\mathrm{J} = \mathrm{J}\)；设 \(p \in \mathrm{J}\)；则 \(p(x^{-1}) = (\lambda(x)p)(e) = 0\)，故 \(x^{-1} \in \mathrm{N}_{\mathrm{J}} = \mathrm{N}\) 且 \(x \in \mathrm{N}\)。这就证明了 (i)。

理想 J 是有限型的（《交换代数》，第三章，§2，第 10 节，定理 2 的推论 2）。因而，存在一个整数 q，使得若 \(W = P_{0} + P_{1} + \cdots + P_{q}\)，则 \(J \cap W\) 作为理想生成 J。用 \(\lambda_{j}\)（相应地，\(\lambda'\)）表示 \(\lambda\) 的由 \(P_{J}\)（相应地，由 W）定义的子表示。由 (1)，

\[
\mathrm{N} = \{x \in \mathbf {G L} _ {n} (k) \mid \lambda^ {\prime} (x) (\mathrm{J} \cap \mathrm{W}) = \mathrm{J} \cap \mathrm{W} \}.\tag{2}
\]

我们证明，对所有 \(j\)，存在表示 \(\sigma_{j}\)（李代数 \(\mathfrak{gl}(n,k)\) 在 \(\mathrm{P}_j\) 上的表示），使得：

\[
\sigma_ {j} | \mathfrak {n} (n, k) \text {   is   compatible   } (\S 3, \text {   no.   } 1) \text {   with   } \lambda_ {j} | T.\tag{3}
\]

\[
\text { For   all } x \in k. 1 _ {n}, \sigma_ {j} (x) \text { is   a   homothety. }\tag{4}
\]

由于 \(\lambda_{j}\) 是 \(\lambda_{1}\) 的第 j 个对称幂，只需证明 \(\sigma_{1}\) 的存在性，参见引理 1。现在 \(\lambda_{1}\) 是表示 \(\gamma\)（\(\mathbf{GL}_{n}(k)\) 在 \(\mathfrak{gl}(n,k)\) 上的表示）的反步表示，由下式给出

\[
\gamma (x) y = x y, \quad x \in \mathbf {G L} _ {n} (k), \quad y \in \mathfrak {g l} (n, k).
\]

设 \(c\) 是李代数 \(\mathfrak{gl}(n,k)\) 在 \(\mathfrak{gl}(n,k)\) 上的表示，由下式给出

\[
c (x) y = x y, \quad x, y \in \mathfrak {g l} (n, k).
\]

直接看出 \(c|\mathfrak{n}(n,k)\) 与 \(\gamma|T\) 相容，且 \(c(x)\) 对所有 \(x \in k.1_{n}\) 是位似。于是，只需取 \(\sigma_{1}\) 为 c 的对偶表示（第一章，§3，第 3 节）。

现在设 \(\sigma'\) 是 \(\mathfrak{gl}(n,k)\) 在 W 上的表示，由诸 \(\sigma_j\)（\(0 \leq j \leq q\)）的直和给出。鉴于 (2) 与关系式

\[
\lambda^ {\prime} (\exp (x)) = \exp (\sigma^ {\prime} (x)) \text { and } \sigma^ {\prime} (\log (y)) = \log (\lambda^ {\prime} (y)), x \in \mathfrak {n} (n, k), y \in \mathrm{T},
\]

我们有

\[
\mathfrak {n} = \{x \in \mathfrak {n} (n, k) \mid \sigma^ {\prime} (x) (\mathrm{J} \cap \mathrm{W}) \subset \mathrm{J} \cap \mathrm{W} \}.\tag{5}
\]

设 \(d = \dim(\mathrm{J} \cap \mathrm{W})\)，并设 \(\tau = \bigwedge^{d} \sigma'\)。设 \(\mathrm{D} = \bigwedge^{d} (\mathrm{J} \cap \mathrm{W})\)。由 (5) 与引理 2 (i)，

\[
\mathfrak {n} = \{x \in \mathfrak {n} (n, k) \mid \tau (x) (\mathrm{D}) \subset \mathrm{D} \}.\tag{6}
\]

但 \(\tau (\mathfrak{n}(n,k))\) 由幂零自同态组成，故 (6) 也可写成

\[
\mathfrak {n} = \{x \in \mathfrak {n} (n, k) \mid \tau (x) (\mathrm{D}) = 0 \}.\tag{7}
\]

现在设 \(E = \bigwedge^{d} W \oplus \bigwedge^{1} V \oplus \bigwedge^{2} V \oplus \cdots \oplus \bigwedge^{n} V\)；设 \(\rho\) 是 \(\tau\) 与 \(\mathfrak{gl}(n, k)\) 在 \(\bigwedge^{1} V, \ldots, \bigwedge^{n} V\) 上的典则表示的直和。设 \(E_{0} \subset E\) 是 \(D = \bigwedge^{d} (J \cap W)\) 与由 \(e_{n-j+1} \wedge \cdots \wedge e_{n}\)（对 \(j = 1, \ldots, n\)）生成的直线的和。由 (7) 与引理 2 (ii)，

\[
\mathfrak {n} = \{x \in \mathfrak {g l} (\mathrm{V}) \mid \rho (x) (\mathrm{E} _ {0}) = 0 \}.\tag{8}
\]

直接看出，若 \(x \in k.1_{n}\)，则 \(\rho(x)\) 可对角化。最后，若 F 是 E 中被 \(\rho(\mathfrak{n})\) 零化的元素所成的集，则 F 在 \(\rho(\mathfrak{q})\) 下稳定（第一章，§3，第 5 节，命题 5），且由 (8)，

\[
\mathfrak {n} = \{x \in \mathfrak {g l} (\mathrm{V}) \mid \rho (x) (\mathrm{F}) = 0 \}.\tag{9}
\]

